\documentclass[11pt,a4paper]{article}
\usepackage[utf8]{inputenc}
\usepackage[T5]{fontenc}
\usepackage[vietnamese]{babel}
\usepackage{lmodern}
\usepackage[margin=2.5cm]{geometry}
\usepackage{amsmath}
\usepackage{booktabs}
\usepackage[hidelinks]{hyperref}
\title{Báo cáo thực nghiệm E2A\\So sánh Representation Pipelines của Frozen DINOv3}
\author{}
\date{15 tháng 9 năm 2026}
\begin{document}
\maketitle

\section{Tóm tắt và quyết định}

E2A so sánh bốn representation pipelines cho fine-grained image retrieval
trên CUB-200-2011 trước khi xây dựng pair-wise confidence learning ở E2B.
Theo protocol đã sửa để chọn phương pháp trên classes 100--199,
\textbf{M1 (final CLS token) là winner}: Hits@1 = 5,214/5,924 và
Recall@1 = 88.0149\%. M3 đạt 5,171 Hits@1, thấp hơn M1 43 queries,
tương đương 0.7259 điểm phần trăm. Không cần tie breaker.
Đây là \textbf{test-based method selection}, không phải kết quả trên
final test độc lập chưa từng được dùng để chọn phương pháp.

\section{Thiết lập và phương pháp}

Backbone frozen DINOv3 ViT-B/16, embedding 768 chiều, checkpoint
\texttt{facebook/dinov3-vitb16-pretrain-lvd1689m}, revision
\texttt{5931719e67bbdb9737e363e781fb0c67687896bc}.
Processor dùng ảnh 224$\times$224 và ImageNet mean/std theo cache manifest.
Có 4 register tokens và 196 final-layer patch tokens; register tokens không
tham gia patch pooling. E2A không dùng EDL, uncertainty hay reranking.

Development gồm 5,864 ảnh classes 0--99, stratified image-level 80/20 với
seed 42: 4,687 fit và 1,177 validation. Test selection gồm 5,924 ảnh thuộc
unseen classes 100--199. Không dùng official image-level split của CUB.
Query và gallery là cùng tập trong mỗi split; self-match loại bằng image ID.
Evaluator dùng L2 normalization, exact cosine similarity, stable
gallery-index tie breaking và Recall@1/2/4/8; lưu Top-100 candidate IDs/scores.

Gọi $c$ là CLS, $p_i$ là patch token và $\bar p$ là patch mean:
\begin{align}
\text{M1:}\quad z &= c,\\
\text{M2:}\quad z &= \bar p=\frac{1}{196}\sum_{i=1}^{196}p_i,\\
\text{M3:}\quad z &= W[c;\bar p]+b,\\
\text{M4:}\quad a_i &= \operatorname{softmax}_i
  \bigl(w_2^\top\tanh(W_1p_i+b_1)+b_2\bigr),\quad z=\sum_i a_ip_i.
\end{align}
M1/M2 không có trainable component. M3 dùng affine projection
$1536\rightarrow768$. M4 dùng Attention Pooling $768\rightarrow256\rightarrow1$,
học content-dependent weights trên final-layer patch-token representations
để tạo global retrieval embedding. Attention weights chỉ dùng diagnostic,
không phải bằng chứng localization hay causal importance.

M3/M4 được fit bằng Proxy Anchor\footnote{S. Kim, D. Kim, M. Cho, S. Kwak,
\href{https://arxiv.org/abs/2003.13911v1}{Proxy Anchor Loss for Deep Metric
Learning}, arXiv:2003.13911v1, CVPR 2020. E2A là adaptation trên frozen
DINOv3, không phải faithful reproduction toàn bộ hệ thống trong bài báo.}
với seed 42, 30 epochs, global batch 20 classes $\times$ 4 images;
AdamW, LR $10^{-4}$, weight decay $10^{-4}$ cho representation và 0 cho
proxies, cosine scheduler, scale 32, margin 0.1, gradient clipping 5.0.
Configs hiện tại dùng FP32 training cho cả M3/M4; backbone extraction có
thể dùng FP16 AMP. Runtime mục tiêu là hai T4, một process/GPU với DDP
và global differentiable gather trước Proxy Anchor.

\begin{table}[ht]
\centering
\caption{Số tham số ngoài frozen backbone. Optimized gồm proxies;
deployable chỉ gồm representation. Epoch lấy từ manifest của run hiện tại.}
\begin{tabular}{lrrr}
\toprule
Method & Optimized & Deployable & Selected epoch\\
\midrule
M1 & 0 & 0 & --\\
M2 & 0 & 0 & --\\
M3 & 1,257,216 & 1,180,416 & 28\\
M4 & 273,921 & 197,121 & 27\\
\bottomrule
\end{tabular}
\end{table}

\section{Protocol lựa chọn}

Validation chọn checkpoint M3/M4 theo integer Hits@1 $\rightarrow$ Hits@2
$\rightarrow$ Hits@4 $\rightarrow$ Hits@8 $\rightarrow$ earliest epoch.
Không dùng validation/test rows để tính training gradient.
Sau validation, chạy cả bốn phương pháp trên classes 100--199 và chọn
winner theo integer test Hits@1 $\rightarrow$ Hits@2 $\rightarrow$ Hits@4
$\rightarrow$ Hits@8 $\rightarrow$ fewer optimized parameters
$\rightarrow$ fixed M1--M4 order. Không còn validation-based selection lock.
Tolerance float32 chỉ dùng để đối chiếu artifacts, không chọn winner.

\section{Kết quả validation và test selection}

\begin{table}[ht]
\centering
\caption{Validation Recall (\%), 1,177 queries. M3 tốt nhất ở validation,
nhưng validation không quyết định winner trong protocol đã sửa.}
\begin{tabular}{lrrrr}
\toprule
Method & R@1 & R@2 & R@4 & R@8\\
\midrule
M1 & 83.52 & 88.95 & 93.54 & 96.43\\
M2 & 46.81 & 57.43 & 70.09 & 80.71\\
M3 & 86.32 & 91.08 & 94.65 & 96.60\\
M4 & 77.06 & 84.79 & 91.08 & 94.82\\
\bottomrule
\end{tabular}
\end{table}

\begin{table}[ht]
\centering
\caption{Test-selection Recall (\%) trên 5,924 queries, từ bundle mới nhất.}
\begin{tabular}{lrrrr}
\toprule
Method & R@1 & R@2 & R@4 & R@8\\
\midrule
\textbf{M1} & \textbf{88.01} & \textbf{92.76} & 95.26 & \textbf{97.01}\\
M2 & 45.76 & 57.87 & 68.23 & 78.75\\
M3 & 87.29 & 92.27 & 95.26 & 96.88\\
M4 & 80.65 & 88.37 & 92.83 & 95.91\\
\bottomrule
\end{tabular}
\end{table}

\begin{table}[ht]
\centering
\caption{Integer test Hits dùng để chọn winner, không dùng Recall làm tròn.}
\begin{tabular}{lrrrr}
\toprule
Method & Hits@1 & Hits@2 & Hits@4 & Hits@8\\
\midrule
M1 & 5,214 & 5,495 & 5,643 & 5,747\\
M2 & 2,711 & 3,428 & 4,042 & 4,665\\
M3 & 5,171 & 5,466 & 5,643 & 5,739\\
M4 & 4,778 & 5,235 & 5,499 & 5,682\\
\bottomrule
\end{tabular}
\end{table}

M3 tăng validation R@1 so với M1 khoảng 2.8037 điểm phần trăm, nhưng
thấp hơn M1 trên unseen-class selection split. Điều này phù hợp với
limitation của same-class validation, nhưng chưa chứng minh nguyên nhân
là overfitting. M2 thấp nhất ở cả hai split. M4 tốt hơn mean pooling,
nhưng chưa vượt CLS baseline. Trong cấu hình này, learned aggregation
không cải thiện test-selection R@1 so với final CLS.

M1 hơn M3 lần lượt 43, 29, 0 và 8 Hits ở K=1/2/4/8. Đây là observed
differences; chưa tính paired confidence interval hay significance test,
do đó không kết luận statistical superiority.

\section{Kiểm tra artifacts và tái lập}

Nguồn chính là \texttt{outputs/e2a\_selection/test\_selection.json} và
\texttt{artifacts/m1/} đến \texttt{m4/}; validation lấy từ
\texttt{metrics/validation.json} ở các thư mục phương pháp.
Đã kiểm tra trực tiếp common 5,924 IDs duy nhất, labels classes 100--199,
Top-100 shape $[5924,100]$, candidates đúng split, không self-match hay
candidate trùng; cosine scores hữu hạn và giảm dần. Recomputed Hits,
Recall và config/metrics/ranking hashes khớp selection JSON.

Test embeddings M3/M4 shape $[5924,768]$, hữu hạn và nonzero. Cosine scores
kiểm tra trên subset 64 queries sai lệch tối đa dưới $1.1\times10^{-6}$;
đây không phải full ranking recomputation. M4 attention weights shape
$[5924,196]$, không âm, tổng gần 1; entropy hữu hạn và khớp weights trong
sai số lưu FP16. M1/M2 dùng manifest-addressed full feature caches
$[11788,768]$, không có test embedding riêng trong portable bundle.
Bundle không kèm backbone caches, patch shards hay training checkpoints;
kiểm tra này không thay thế tái chạy toàn bộ training/extraction.

\section{Limitations và bàn giao E2B}

\begin{itemize}
\item Classes 100--199 đã được dùng để chọn phương pháp và đã có rerun;
không còn untouched final test. E2A/E2B dùng lại split này phải disclosure
selection bias, không gọi là unbiased final generalization.
\item Fit/validation khác ảnh nhưng cùng classes 0--99; validation chỉ
đo image generalization trong cùng class set, không trực tiếp đo
unseen-class generalization.
\item M3/M4 là single-seed results (42); không báo mean $\pm$ std theo
seed, chưa đo training variance. M1/M2 không có training-seed variance.
\item Đây là comparison of representation pipelines, không phải ablation
chỉ thay pooling: learned components và số tham số khác nhau.
\item FP16 extraction/storage có thể tạo thay đổi nhỏ giữa reruns.
Báo cáo dùng bundle mới nhất, không trộn kết quả của các lần chạy.
Attention diagnostics không phải localization evidence.
\end{itemize}

Quyết định bàn giao: dùng \textbf{M1 final CLS}, L2 và cosine làm
representation/R0 cho E2B, giữ nguyên split, preprocessing, self-match
exclusion và Top-100 contract. Confidence learning vẫn phải fit/tune trên
development data. Báo cáo E2B cần kế thừa disclosure rằng representation
được chọn bằng test performance.

\end{document}
